\chapter{Information-Form Identities}

This appendix records supporting coordinate and factor-assembly identities that are used by the main construction. It follows the main text so that these details support, rather than precede, the definitions and ELBO dependency order.

\section{Gaussian completion, normalization, and dual derivatives}

Fix $h\in\mathbb R^D$ and $J=J^\top\succ0$, and define $\mu:=J^{-1}h$.  Direct expansion gives
\begin{equation}
\begin{aligned}
h^\top y-\frac12y^\top Jy
&=-\frac12
\left(y-J^{-1}h\right)^\top
J
\left(y-J^{-1}h\right)
+\frac12h^\top J^{-1}h\\
&=-\frac12(y-\mu)^\top J(y-\mu)
+\frac12h^\top J^{-1}h.
\end{aligned}
\label{eq:information-form-completion-square}
\end{equation}
The positive-definite Gaussian integral is
\begin{equation}
\int_{\mathbb R^D}
\exp\left[-\frac12(y-\mu)^\top J(y-\mu)\right]
\mathop{\mathrm dy}
=(2\pi)^{D/2}(\det J)^{-1/2}.
\label{eq:precision-gaussian-integral}
\end{equation}
Combining \cref{eq:information-form-completion-square,eq:precision-gaussian-integral} yields
\begin{equation}
\log\int_{\mathbb R^D}
\exp\left(h^\top y-\frac12y^\top Jy\right)
\mathop{\mathrm dy}
=\frac12h^\top J^{-1}h
-\frac12\log\det J
+\frac D2\log(2\pi),
\label{eq:derived-gaussian-log-normalizer}
\end{equation}
which proves \cref{eq:gaussian-log-normalizer}.  The completed density is $\mathcal N(\mu,J^{-1})$, establishing \cref{eq:information-moment-conversion}.

For the derivative calculation, set $\Sigma:=J^{-1}$ and $M:=\Sigma+\mu\mu^\top$.  The inverse and determinant differentials are
\begin{equation}
\mathop{\mathrm d}(J^{-1})
=-J^{-1}(\mathop{\mathrm d}J)J^{-1},
\qquad
\mathop{\mathrm d}\log\det J
=\operatorname{tr}(J^{-1}\mathop{\mathrm d}J).
\label{eq:inverse-logdet-differentials}
\end{equation}
Holding the other argument fixed for each partial differential gives
\begin{equation}
\begin{aligned}
\mathop{\mathrm d}_hA
&=\mu^\top\mathop{\mathrm d}h,\\
\mathop{\mathrm d}_JA
&=-\frac12\mu^\top(\mathop{\mathrm d}J)\mu
-\frac12\operatorname{tr}(\Sigma\mathop{\mathrm d}J)\\
&=-\frac12\langle M,\mathop{\mathrm d}J\rangle_{\mathrm F}.
\end{aligned}
\label{eq:log-normalizer-differential-proof}
\end{equation}
Because $M$ and $J$ are symmetric, the Frobenius gradients are exactly $\nabla_hA=\mu$ and $\nabla_JA=-M/2$.  Under the linear reparameterization $\eta_2=-J/2$, the chain rule gives $\nabla_{\eta_2}A=M$, confirming that $(\mu,M)$, rather than $(\mu,\Sigma)$, is the expectation-coordinate partner of $(h,-J/2)$.

\section{Joint entropy and Gaussian KL derivations}

Taking expectation in \cref{eq:joint-gaussian-information-form} and using
\begin{equation}
\mathbb E_q[y^\top Jy]
=\operatorname{tr}(JM)
=D+\mu^\top J\mu,
\qquad
h^\top\mu=\mu^\top J\mu,
\label{eq:gaussian-quadratic-expectation}
\end{equation}
gives
\begin{equation}
\begin{aligned}
H(q)
&=A(h,J)-h^\top\mu
+\frac12\operatorname{tr}(JM)\\
&=\frac12
\left[D\bigl(1+\log(2\pi)\bigr)-\log\det J\right],
\end{aligned}
\label{eq:joint-entropy-derivation}
\end{equation}
including the full joint determinant.  If $q_t$ are the token marginals, then
\begin{equation}
\begin{aligned}
D_{\mathrm{KL}}
\left(q(y)\middle\Vert\prod_tq_t(y_t)\right)
&=\mathbb E_q\left[
\log q(y)-\sum_t\log q_t(y_t)
\right]\\
&=\sum_tH(q_t)-H(q),
\end{aligned}
\label{eq:total-correlation-entropy-proof}
\end{equation}
which proves the sign in \cref{eq:gaussian-total-correlation}.

For $q_i=\mathcal N(\mu_i,\Sigma_i)$, write $\delta:=\mu_0-\mu_1$.  The two quadratic expectations needed for the oriented ratio $\log(q_0/q_1)$ are
\begin{equation}
\begin{aligned}
\mathbb E_{q_0}
[(y-\mu_0)^\top\Sigma_0^{-1}(y-\mu_0)]
&=D,\\
\mathbb E_{q_0}
[(y-\mu_1)^\top\Sigma_1^{-1}(y-\mu_1)]
&=\operatorname{tr}(\Sigma_1^{-1}\Sigma_0)
+\delta^\top\Sigma_1^{-1}\delta.
\end{aligned}
\label{eq:gaussian-kl-quadratic-expectations}
\end{equation}
Substitution into $\mathbb E_{q_0}[\log q_0-\log q_1]$ proves \cref{eq:gaussian-kl-covariance-form}.  Replacing $\Sigma_i$ by $J_i^{-1}$ and using
\begin{equation}
\log\frac{\det\Sigma_1}{\det\Sigma_0}
=\log\frac{\det J_0}{\det J_1}
\label{eq:gaussian-kl-determinant-conversion}
\end{equation}
proves \cref{eq:gaussian-kl-precision-form} with the same orientation.  Finally, the exponential-family log ratio gives
\begin{equation}
\begin{aligned}
D_{\mathrm{KL}}(q_0\Vert q_1)
={}&A_1-A_0
-(h_1-h_0)^\top\mu_0\\
&+\frac12\langle J_1-J_0,M_0\rangle_{\mathrm F},
\end{aligned}
\label{eq:natural-expectation-kl-derivation}
\end{equation}
which is \cref{eq:gaussian-kl-natural-expectation-form}.

\section{Gauge pushforward and precision-family closure}

Let $y'=\mathsf G y$ with the block-diagonal invertible map in \cref{eq:block-diagonal-global-frame-map}.  Substituting $y=\mathsf G^{-1}y'$ into the exponent gives
\begin{equation}
h^\top y-\frac12y^\top Jy
=(\mathsf G^{-\top}h)^\top y'
-\frac12y'^\top
(\mathsf G^{-\top}J\mathsf G^{-1})y',
\label{eq:gauge-transformed-information-exponent}
\end{equation}
which proves the $(h,J)$ transformations in \cref{eq:joint-gaussian-gauge-laws}.  Inversion and expectation then give
\begin{equation}
\begin{aligned}
(J')^{-1}
&=\mathsf GJ^{-1}\mathsf G^\top,\\
\mu'=(J')^{-1}h'
&=\mathsf G\mu,\\
M'=\mathbb E[y'y'^\top]
&=\mathsf GM\mathsf G^\top.
\end{aligned}
\label{eq:gauge-transformed-moment-proof}
\end{equation}
Moreover,
\begin{equation}
\log\det J'
=\log\det J-2\log|\det\mathsf G|,
\label{eq:gauge-precision-logdet-shift}
\end{equation}
while the quadratic term $h^\top J^{-1}h$ is unchanged.  Hence $A'=A+\log|\det\mathsf G|$ and $H'=H+\log|\det\mathsf G|$.  The pushed-forward density acquires the inverse Jacobian, so both factors in a KL ratio acquire the same additive log-density shift.  Their difference is unchanged, proving \cref{eq:entropy-shift-kl-invariance}.

\paragraph{Proposition (closure of token-block precision sparsity).}
Let $\mathcal S$ be any set of token-index pairs and let
\begin{equation}
\mathcal J_{\mathcal S}
:=\{J=J^\top\succ0:J_{ts}=0\text{ for }(t,s)\notin\mathcal S\}.
\label{eq:token-block-sparse-precision-family}
\end{equation}
For every block-diagonal $\mathsf G=\operatorname{diag}(\mathsf G_0,\ldots,\mathsf G_T)$ with invertible blocks, the congruence $J'=\mathsf G^{-\top}J\mathsf G^{-1}$ maps $\mathcal J_{\mathcal S}$ onto itself.

\emph{Proof.}
Block multiplication gives $J_{ts}'=\mathsf G_t^{-\top}J_{ts}\mathsf G_s^{-1}$.  Thus $J_{ts}=0$ implies $J_{ts}'=0$.  Applying the same argument to $J=\mathsf G^\top J'\mathsf G$ proves the converse.  Congruence by an invertible matrix preserves positive definiteness, so every transformed matrix remains in the precision cone.  \hfill$\square$

The conclusion does not extend to elementwise diagonality within a token block.  The matrix in \cref{eq:diagonal-covariance-nonclosure-example} has determinant one, yet its congruence image of $I_2$ has off-diagonal entry one.  General $\mathrm{GL}(2)$ therefore takes the diagonal-covariance family outside itself.  The same example applies to a diagonal precision by replacing the frame map with its inverse transpose.

Finally, direct multiplication verifies that the two matrices in \cref{eq:tridiagonal-precision-dense-covariance} are inverses.  Hence even the shortest nondegenerate Markov chain can have a tridiagonal precision and a dense covariance.  Sparse storage promises exact factorization and solve interfaces for $J$; it does not promise sparse storage for $J^{-1}$.

\chapter{Proof Details and Verification Oracles}

This appendix collects proof details and reference checks for claims introduced in the main text. It will make the corresponding assumptions and numerical oracles explicit when those derivations are added.

\section{Full normalization proof for the fixed conditional joint}

By \cref{eq:nonempty-causal-parent-sets}, every state or model source index lies on a permitted edge $(j,t)\in E_T$.  The transports $\Omega^z_{tj}$ and $\Omega^m_{tj}$ used by the corresponding Gaussian kernels are therefore defined by \cref{eq:typed-associated-transports}.

Let $\#_A$ denote counting measure on a finite set $A$.  In the local frames used to express the densities, the reference measure for the fixed-horizon variables $(x,\boldsymbol z,\boldsymbol m,a,b)$ is
\begin{equation}
\begin{aligned}
\mu_T
:={}&\left(\bigotimes_{t=1}^{T}\#_{\mathcal X}\right)
\otimes\left(\bigotimes_{t=0}^{T}
(\lambda_{z,t}\otimes\lambda_{m,t})\right)\\
&\otimes\left(\bigotimes_{t=1}^{T}
(\#_{\operatorname{Pa}_z(t)}\otimes
\#_{\operatorname{Pa}_m(t)})\right).
\end{aligned}
\label{eq:full-joint-reference-measure}
\end{equation}
Each factor is measurable by the assumptions preceding \cref{eq:vfe4-generative-joint}.  The finite products are therefore measurable, and the density is nonnegative, so Tonelli's theorem permits every integration and summation below without a prior integrability assumption.

For fixed earlier variables and fixed $\Gamma$, collect the factors at time $t$ into
\begin{equation}
\begin{aligned}
F_t
:={}&\pi_t^m(b_t)
K^m_{\theta,tb_t}(m_t\given m_{b_t},x_{<t},\Gamma)
\pi_t^z(a_t)\\
&\times K^z_{\theta,ta_t}(z_t\given z_{a_t},m_t,x_{<t},\Gamma)
L_{\theta,t}(x_t\given z_t,m_t,\Gamma).
\end{aligned}
\label{eq:time-slice-factor-for-normalization}
\end{equation}
Nested normalization gives the pointwise identity
\begin{equation}
\begin{aligned}
&\sum_{b_t\in\operatorname{Pa}_m(t)}\pi_t^m(b_t)
\int_{E_{m,t}}K^m_{\theta,tb_t}
\sum_{a_t\in\operatorname{Pa}_z(t)}\pi_t^z(a_t)
\int_{E_{z,t}}K^z_{\theta,ta_t}
\sum_{x_t\in\mathcal X}L_{\theta,t}(x_t\given z_t,m_t,\Gamma)
\mathop{\mathrm d\lambda_{z,t}(z_t)}
\mathop{\mathrm d\lambda_{m,t}(m_t)}\\
&\quad=
\sum_{b_t\in\operatorname{Pa}_m(t)}\pi_t^m(b_t)
\int_{E_{m,t}}K^m_{\theta,tb_t}
\sum_{a_t\in\operatorname{Pa}_z(t)}\pi_t^z(a_t)
\mathop{\mathrm d\lambda_{m,t}(m_t)}\\
&\quad=
\sum_{b_t\in\operatorname{Pa}_m(t)}\pi_t^m(b_t)
\int_{E_{m,t}}K^m_{\theta,tb_t}
\mathop{\mathrm d\lambda_{m,t}(m_t)}
=1.
\end{aligned}
\label{eq:time-slice-nested-normalization}
\end{equation}
The unshown conditioning arguments of the two transition kernels are exactly those in \cref{eq:time-slice-factor-for-normalization}.  In the first equality, categorical emission normalization removes the innermost sum, state-kernel normalization removes the $z_t$ integral, and \cref{eq:normalized-causal-source-priors} removes the $a_t$ sum.  In the second equality, model-kernel normalization removes the $m_t$ integral and the same source-prior identity removes the $b_t$ sum.

To apply \cref{eq:time-slice-nested-normalization} to the complete density, begin at the terminal vertex $T$.  There is no factor at an earlier time that depends on $x_T$, $z_T$, $m_T$, $a_T$, or $b_T$, because every parent index is strictly smaller than its receiver and every history is $x_{<t}$.  Factors at later times may depend on terminal-block variables at an intermediate stage, but reverse topological order removes all such factors before that stage is reached.  Hence the nested elimination at $T$ is pointwise in every earlier variable and replaces $F_T$ by one.

After eliminating the terminal block, the same argument applies to $F_{T-1}$.  Induction removes $F_t$ for $t=T,T-1,\ldots,1$.  History dependence is harmless throughout: when the time-$t$ identity is evaluated, $x_{<t}$ is held fixed, and all factors that could depend on $x_t$ have already been eliminated.  What remains is
\begin{equation}
\int_{E_{z,0}\times E_{m,0}}
p_{\theta,0}(z_0,m_0\given\Gamma)
\mathop{\mathrm d\lambda_{z,0}(z_0)}
\mathop{\mathrm d\lambda_{m,0}(m_0)}
=1,
\label{eq:normalization-proof-initial-remainder}
\end{equation}
by \cref{eq:normalized-initial-law}.  Thus the density over $(x,\boldsymbol z,\boldsymbol m,a,b)$ in \cref{eq:vfe4-generative-joint} defines a probability measure with respect to \cref{eq:full-joint-reference-measure} for every fixed deterministic $\Gamma$.  Under a local frame change, \cref{eq:receiver-density-kernel-measure-invariance,eq:initial-density-gauge-jacobian} preserve this probability measure even though its coordinate density and Lebesgue representatives change.  \hfill$\square$

\section{Bundle-map dimension and type oracle}

The coordinate audit for \cref{eq:typed-associated-transports,eq:typed-model-state-morphism} begins with column vectors $z_j,z_t\in\mathbb R^{d_z}$ and $m_j,m_t\in\mathbb R^{d_m}$.  The represented state link has shape $d_z\times d_z$, the represented model link has shape $d_m\times d_m$, and the cross-channel morphism has shape $d_z\times d_m$.  Consequently,
\begin{equation}
\begin{aligned}
\underbrace{\Omega^z_{tj}}_{d_z\times d_z}
\underbrace{z_j}_{d_z\times 1}
&\in\mathbb R^{d_z},
&
\underbrace{\Omega^m_{tj}}_{d_m\times d_m}
\underbrace{m_j}_{d_m\times 1}
&\in\mathbb R^{d_m},\\
\underbrace{B_t}_{d_z\times d_m}
\underbrace{m_t}_{d_m\times 1}
&\in\mathbb R^{d_z}.
&&
\end{aligned}
\label{eq:bundle-map-dimension-oracle}
\end{equation}
The first and third outputs in \cref{eq:bundle-map-dimension-oracle} inhabit the same target fiber $E_{z,t}$ and may enter one state-kernel location parameter.  The model-link output inhabits $E_{m,t}$ and may enter a model-kernel location parameter.  Combining the latter directly with a state-fiber vector would be ill typed.

The gauge-law audit is equally mechanical.  With $\rho_z(g_t)\in\mathbb R^{d_z\times d_z}$ and $\rho_m(g_t)\in\mathbb R^{d_m\times d_m}$, every product in \cref{eq:link-morphism-gauge-laws} has the same shape as the unprimed map:
\begin{equation}
\begin{aligned}
(d_z\times d_z)(d_z\times d_z)(d_z\times d_z)&=d_z\times d_z,\\
(d_m\times d_m)(d_m\times d_m)(d_m\times d_m)&=d_m\times d_m,\\
(d_z\times d_z)(d_z\times d_m)(d_m\times d_m)&=d_z\times d_m.
\end{aligned}
\label{eq:gauge-law-dimension-oracle}
\end{equation}
No equality $d_z=d_m$ is used.  This check is the source-level oracle for later probability kernels.  A covariance on the state fiber is a positive symmetric contravariant tensor and its precision is a positive symmetric covariant tensor,
\begin{equation}
\Sigma_{z,t}\in\operatorname{Sym}^{2}_{+}(E_{z,t})
\simeq\operatorname{Hom}_{+}(E_{z,t}^{*},E_{z,t}),
\qquad
J_{z,t}\in\operatorname{Sym}^{2}_{+}(E_{z,t}^{*})
\simeq\operatorname{Hom}_{+}(E_{z,t},E_{z,t}^{*}).
\label{eq:state-covariance-precision-types}
\end{equation}
The model-fiber tensors have the analogous types with $E_{z,t}$ replaced by $E_{m,t}$.  If
\begin{equation}
\mathsf G_{z,t}:=\rho_z(g_t),
\qquad
\mathsf G_{m,t}:=\rho_m(g_t),
\label{eq:fiber-coordinate-change-matrices}
\end{equation}
then covariance and precision carry opposite congruence laws in either channel:
\begin{equation}
\begin{aligned}
\Sigma_{z,t}'&=\mathsf G_{z,t}\Sigma_{z,t}\mathsf G_{z,t}^{\top},
&J_{z,t}'&=\mathsf G_{z,t}^{-\top}J_{z,t}\mathsf G_{z,t}^{-1},\\
\Sigma_{m,t}'&=\mathsf G_{m,t}\Sigma_{m,t}\mathsf G_{m,t}^{\top},
&J_{m,t}'&=\mathsf G_{m,t}^{-\top}J_{m,t}\mathsf G_{m,t}^{-1}.
\end{aligned}
\label{eq:covariance-precision-opposite-congruence}
\end{equation}
A kernel over $E_{z,t}$ must therefore receive a location in $E_{z,t}$ together with a covariance in $\operatorname{Sym}^{2}_{+}(E_{z,t})$ or a precision in $\operatorname{Sym}^{2}_{+}(E_{z,t}^{*})$.  A kernel over $E_{m,t}$ must receive the corresponding model-fiber objects.

\section{Expanded Regime I composition check}

For a path $v_0\to v_1\to\cdots\to v_r$ in the auxiliary geometric graph, the composite group transport is ordered so that the rightmost map acts first.  Repeated use of \cref{eq:regime-i-cocycle-identity} gives
\begin{equation}
\Omega_{v_rv_{r-1}}\cdots\Omega_{v_1v_0}
=U_{v_r}U_{v_0}^{-1}.
\label{eq:regime-i-open-path-telescope}
\end{equation}
Thus an open-path transport depends only on its endpoints in Regime I.  Setting $v_r=v_0$ yields the identity, as in \cref{eq:regime-i-loop-holonomy}.  Applying either representation to \cref{eq:regime-i-open-path-telescope} gives the correctly typed composites
\begin{equation}
\Omega^z_{v_rv_{r-1}}\circ\cdots\circ\Omega^z_{v_1v_0}
:E_{z,v_0}\longrightarrow E_{z,v_r},
\qquad
\Omega^m_{v_rv_{r-1}}\circ\cdots\circ\Omega^m_{v_1v_0}
:E_{m,v_0}\longrightarrow E_{m,v_r}.
\label{eq:associated-open-path-types}
\end{equation}
This check uses a geometric walk in the underlying undirected graph or cell complex.  It does not insert a directed cycle into the causal factor graph.

\chapter{Notation}

This appendix reserves a stable location for the notation shared by the preceding construction. It does not introduce additional model assumptions beyond those defined in the main text.
